%!TEX root = ../main.tex

\chapter{Fase 6: versione finale consegnabile}

\section{Obiettivo della versione finale}

La Fase 6 produce una relazione finale autonoma, pronta per essere usata in Prism oppure impaginata in PDF, Word o LaTeX. A differenza dei capitoli preparatori, la versione finale non è organizzata per fasi: presenta direttamente il tirocinio, il progetto \builder, le attività svolte, la supervisione della tutor e le competenze acquisite.

Il documento finale deve mantenere il tono di una relazione universitaria di tirocinio. Deve quindi descrivere il lavoro tecnico senza trasformarlo in manuale utente, valorizzando allo stesso tempo il percorso formativo, il ruolo della Prof.ssa Alessandra Lumini e il contributo del progetto alla comprensione della progettazione di basi di dati.

\section{Struttura finale adottata}

\noindent\begin{tabularx}{\textwidth}{>{\bfseries}p{4cm}X}
Sezione finale & Contenuto \\
\hline
Introduzione & Presentazione del tirocinio, del progetto \builder, della tutor e dello scopo della relazione. \\
Obiettivi del tirocinio & Obiettivo generale, obiettivi specifici, valore didattico, correttezza e usabilità. \\
Il progetto \builder & Descrizione dell'applicazione, destinatari, flussi principali e utilità in ambito universitario. \\
Analisi dei requisiti & Requisiti funzionali e non funzionali, con attenzione a modellazione, traduzione, export e qualità. \\
Tecnologie utilizzate & Ruolo di TypeScript, React, Vite, CSS, SVG, npm, Playwright, GitHub e strumenti collegati. \\
Architettura e progettazione & Organizzazione del codice, separazione tra UI e logica, tipi condivisi, testabilità e persistenza. \\
Funzionalità sviluppate e consolidate & Editor ER, traduzione logica, export, salvataggio, interfaccia e localizzazione. \\
Test, verifica e miglioramenti & Test unitari, integrazione, E2E, bug fixing e miglioramenti di stabilità. \\
Supervisione della tutor e competenze acquisite & Ruolo della Prof.ssa Alessandra Lumini, feedback, competenze tecniche e trasversali. \\
Conclusioni e sviluppi futuri & Risultati, limiti realistici e possibili evoluzioni del progetto. \\
\end{tabularx}

\section{Materiali prodotti}

I materiali prodotti in questa fase sono:

\begin{itemize}
    \item una relazione finale autonoma in \texttt{doc/relazione-finale-buildER.tex};
    \item un prompt finale per Prism;
    \item un elenco dei placeholder da completare;
    \item note di compilazione e controlli LaTeX;
    \item suggerimenti sui punti in cui inserire eventuali screenshot o allegati.
\end{itemize}

\section{Indicazioni per la consegna}

Prima della consegna è necessario completare i placeholder, inserire eventuali screenshot, controllare la copertina, verificare l'indice, controllare bibliografia e sitografia, correggere i dati personali e generare il PDF finale. È inoltre consigliata una lettura completa del documento, per verificare tono, coerenza e assenza di ripetizioni.

Gli screenshot più utili riguardano l'interfaccia principale, uno schema ER, la vista logica e un esempio di esportazione. Non sono obbligatori se Prism richiede solo testo, ma possono migliorare la chiarezza della relazione se il formato finale li consente.

\section{Controllo finale}

\begin{itemize}
    \item[$\square$] Copertina completa.
    \item[$\square$] Tutor indicata correttamente.
    \item[$\square$] Corpo centrale circa dieci pagine.
    \item[$\square$] Testo coerente.
    \item[$\square$] Placeholder completati.
    \item[$\square$] Screenshot inseriti, se necessari.
    \item[$\square$] Bibliografia/sitografia controllata.
    \item[$\square$] PDF generato correttamente.
    \item[$\square$] Nessun dato inventato.
    \item[$\square$] Relazione pronta per Prism o consegna.
\end{itemize}
